\appendixtitles{yes}
\appendixstart
\appendix
\def\theHequation{\thesection.\arabic{equation}}\def\theHtheorem{\thesection.\arabic{theorem}}
\section{Energy Balance: Evaluation and Scope}\label{app:balance}

\rev{\emph{Quadrature.} Theorem~\ref{thm:balance} is exact for every $t$; its terms are evaluated by quadrature. The residual $r(\cdot,t)$ is a polynomial of degree $\le3K$ in $x$, so $\|r\|^2$ has degree $\le6K$, and the energy and the terms of the first line of~\eqref{eq:balance-main} have degree $\le4K$. In our implementation the energy and these terms are integrated exactly by Clenshaw--Curtis quadrature on $4K+9$ nodes; $\|r\|^2$ is integrated by the same rule, and its relative deviation from the exact value (Clenshaw--Curtis on $6K+9$ nodes) is at most of the order of \N{quad.max_rel} for the seed-42 network of each configuration of Table~\ref{tab:balance}. The boundary values $u_x(\pm1,t)$ and $u_t(\pm1,t)$ are evaluated directly from the coefficients.}

\rev{\emph{Check of signs and constants.} Identity~\eqref{eq:balance-main} was compared with a central finite difference of $E(t)$, evaluated with exact Gauss--Legendre quadrature, for a random smooth coefficient trajectory with $K=7$ and $\varepsilon=0.3$ at three times: the relative difference was below $10^{-9}$, which is the accuracy of the finite difference (script \texttt{code/check\_identity.py}). This is a plausibility check, not part of the proof.}

\rev{\emph{What the theorem does and does not say.} (i)~It is conditional: it converts a bound on the PDE residual and on the Neumann defect into a bound on energy growth. It does not prove that training makes them small, and the resulting bound $\frac12\int\|r\|^2\,dt+\int\beta\,dt$ can be far larger than the actual energy variation; the time integrals must be measured for each trained network (Table~\ref{tab:balance}). Loss values during training are proxies for them and must not be substituted. (ii)~In the revised Spectral-PINN the PDE loss is the Clenshaw--Curtis quadrature of $\frac12\|r\|^2$ on $N+1$ Gauss--Lobatto nodes, $N=\max(300,4K+8)$, averaged over $100\,T$ times ($250$ for B3), and the Neumann condition holds identically, so $\beta=0$. In the variant of the first version the residual is a plain mean over $300$ equispaced points (not an exact quadrature of the degree-$6K$ polynomial $r^2$) and the boundary penalty acts on $u_x(\pm1,t)^2$ with weight $K^{-4}\lambda_{\rm BC}$, a very weak control of $\beta$ at large $K$; this is the cause of the loss of accuracy at large $K$ reported in the first version (Section~\ref{sec:results-k}). (iii)~The bound does not involve the energy-penalty term of the loss: a small residual and a small Neumann defect already constrain the energy growth, which is consistent with the penalty being inactive (Section~\ref{sec:results-energy}). (iv)~With the initial condition $x^2\cos(\pi x)$ of the first version, $u_{0,x}(-1)=2\neq0$, so for a soft-variant network that fits $u_0$ the factor $u_x(-1,0)$ in $\beta$ is not small at $t=0$; with the compatible initial condition~\eqref{eq:IC} and the exact Neumann basis $\beta=0$ by construction.}

\section{Symmetry-Preservation Diagnostic}\label{app:sym}

For the two-dimensional experiment of Section~\ref{sec:2d}, the reflection-symmetry violation $S(t)$ was evaluated on an endpoint-inclusive grid $x_i=-\pi+ih$, $h=2\pi/(n-1)$, $i=0,\dots,n-1$ ($n=49$), so that $x_{n-1-i}=-x_i$ \rev{up to floating-point round-off (for the computed nodes $|x_{n-1-i}+x_i|\le9\times10^{-16}$)}; this avoids an artificial, purely discretization-induced violation that would otherwise arise from comparing a periodic (endpoint-excluded) training grid against its naive array reversal. \rev{The remaining round-off of the nodes, not the network, produces the nonzero values of order $10^{-31}$ of $S$ for the Fourier-Spectral-PINN in Tables~\ref{tab:2d} and~\ref{tab:stress-2d}: on an exactly mirrored grid the cosine basis gives $S=0$, whereas on the computed grid the field $0.4\cos(2x)\cos(2y)$ by itself already gives $S=8.7\times10^{-32}$.}

\section{Original Initial Condition}\label{app:orig}

\rev{The first version used $u_0(x)=x^2\cos(\pi x)$, which does not satisfy the Neumann condition ($u_0'(\mp1)=\pm2$). For $t>0$ the solution of~\eqref{eq:AC} with this initial condition satisfies the Neumann condition and has a thin initial boundary layer; the certified reference of Appendix~\ref{app:ref} resolves it (the wall unknowns are eliminated with the Neumann rows, so the boundary condition holds exactly for $t>0$). Table~\ref{tab:orig} repeats the B1 experiments of the as-submitted Spectral-PINN with this condition (seeds 42, 123 and 7) and compares them with the compatible condition~\eqref{eq:IC} for the same seeds; each run is scored against the certified reference for its own initial condition.}

\rev{The conclusions that can be tested with the original condition do not depend on the choice. With $K=32$ the errors are \N{orig.K32} (original) and \N{orig.K32c} (compatible). With $K=64$ the as-submitted network fails with both initial conditions (\N{orig.K64} and \N{orig.K64c}), with wall slopes up to \N{orig.K64.wall} and \N{orig.K64c.wall}, so the boundary-penalty artifact of Section~\ref{sec:results-k} is not caused by the incompatible initial condition. The classical solver and the reaction-only model have errors of the same size for both conditions, and all networks have a monotonically decreasing energy (\N{orig.viol} increases on the grid of spacing $0.02$). The exact Neumann basis of the revised network cannot represent an initial condition with a nonzero wall slope, which is a further reason for using a compatible one in the main text.}

\begin{table}[H]
\caption{\rev{Original initial condition $x^2\cos(\pi x)$ versus the compatible condition~\eqref{eq:IC}: $L^2$ error at $t=1$ ($\varepsilon=0.01$) of the as-submitted Spectral-PINN, mean $\pm$ SD over $n$ seeds, and the largest wall slope $|u_x(\pm1,t)|$.}\label{tab:orig}}
\small
\begin{tabularx}{\textwidth}{Xcccccc}
\toprule
Variant & $n$ & original: $L^2$ & wall slope & $n$ & compatible: $L^2$ & wall slope\\
\midrule
as submitted, $K=32$ & 0 & \textbf{--} & \textbf{--} & 0 & \textbf{--} & \textbf{--}\\
as submitted, $K=64$ & 0 & \textbf{--} & \textbf{--} & 4 & \ensuremath{(9.93\pm1.56)\times10^{-2}} & 5.67\\
\bottomrule
\end{tabularx}

\end{table}


\section{Certified Reference Solution}\label{app:ref}

\rev{All accuracy figures of this article are measured against a spectral reference solution computed as follows. The unknowns are the values of $u$ at the Chebyshev--Gauss--Lobatto nodes $x_j=\cos(j\pi/N)$; the two rows $(Du)_0=(Du)_N=0$ of the differentiation matrix $D$ are used to eliminate the wall unknowns (equivalent to a differential--algebraic system with algebraic boundary rows), so that the Neumann condition holds exactly for $t>0$ also when $u_0$ does not satisfy it. Both initial conditions used here are even in $x$, and the equation and the Neumann condition commute with $x\mapsto-x$, so the solver works on the even half ($N/2$ unknowns), which was verified to agree with the full-interval solver to $2\times10^{-13}$. The time integrator is the fifth-order L-stable Radau method (adaptive, analytic Jacobian, rtol $10^{-11}$, atol $10^{-13}$), restarted at every output time so that the returned values are step endpoints. Values at the evaluation points are obtained by barycentric interpolation. The resolution is $N=128$ for $\varepsilon\ge0.05$ and $N=256$ for $\varepsilon=0.01$ ($T\le1$), and $N=\min(2048,\max(256,2\lceil5.12/\varepsilon\rceil))$ for $T>1$.}

\rev{\emph{Certification.} ``Certified'' means demonstrated numerically, not a rigorous a-posteriori bound. (i)~For $T\le1$, $\varepsilon\in\{0.01,0.05,0.1\}$ and both initial conditions, the maximum difference of the default solution at $t\in\{0.1,0.5,1\}$ from the solution with $N=768$ is $\le4\times\nobreak10^{-11}$ (the floor of $\approx3.5\times10^{-11}$ is rounding in the large differentiation matrices); tightening rtol to $10^{-12}$ changes the solution by $\le4\times10^{-13}$ (tested at $N=384$); the BDF and LSODA integrators agree to $\le1.3\times10^{-10}$. (ii)~An independent second-order finite-difference discretization with ghost nodes for the Neumann condition ($N_x=1000,2000,4000,8000$, Radau) converges at the rate $h^2$ and, after Richardson extrapolation, agrees with the spectral solution to $\le5.4\times10^{-9}$ (worst case: the incompatible initial condition at $\varepsilon=0.01$, $t=0.1$) and mostly to $10^{-12}$--$10^{-10}$; since the two discretizations are coded separately, this rules out an error in the sign, scaling or boundary rows of the spectral solver. (iii)~For the interface benchmark B3 ($\varepsilon=0.05$, $T=5$) the default $N=256$ differs from $N=384$, $512$ and $768$ by $2.7\times10^{-12}$, $6.4\times10^{-13}$ and $2\times10^{-11}$ in the maximum norm over $t\in\{0.2,0.5,1,2,3,4,5\}$. (iv)~For $\varepsilon=0.01$ and $T=10$ one needs $N\ge1024$ ($N=256$ is wrong by $10^{-2}$); this case is not used here.}

\rev{\emph{The reference of the first version.} The first version used a semi-implicit finite-difference scheme (first order in time, second order in the interior; $N_x=500$, $\delta t=10^{-4}$) with the wall rows $(1+\lambda)u_0-\lambda u_1$, which are first-order accurate and equivalent to halving the diffusion at the two wall nodes; the ghost-node rows $(1+2\lambda)u_0-2\lambda u_1$ are the consistent ones. Against the certified solution its maximum error on the 200-point grid was $2.7$--$3.7\times10^{-3}$ at $t=0.1$, $1.4$--$2.5\times10^{-3}$ at $t=0.5$ and $5.6\times10^{-4}$--$1.4\times10^{-3}$ at $t=1$ for $\varepsilon\in\{0.01,0.05,0.1\}$ (original initial condition). Re-scoring the stored predictions of the first version against the certified solution changed the $L^2$ errors of the Spectral-PINN by at most $2.2\%$ and those of the vanilla PINN by at most $0.3\%$ (the maximum-norm errors by at most $0.5\%$ and $5.2\%$), so the PINN rows were not materially affected; the classical baseline row was (its error of $8.5\times10^{-5}$ was mostly the error of the reference, and the norm used for that row was inconsistent with the PINN rows: the sum over the $N+1$ Chebyshev nodes was multiplied by the spacing of the finite-difference grid instead of integrating over $[-1,1]$).}

\rev{\emph{Reaction-only comparison.} The closed-form solution of $u'=u-u^3$ at each $x$, $u(x,t)=u_0e^t/\sqrt{1+u_0^2(e^{2t}-1)}$, is used as a reference for ``how much of the solution is explained by the reaction term alone'': its distance to the certified solution is the number given as ``Reaction-only ODE ($\varepsilon=0$)'' in Tables~\ref{tab:forward} and~\ref{tab:orig}.}

